\newcommand{\AbstractResult}{三个随机种子的平均结果表明，雷达补检将纯相机的 mAP50--95 从 1.38\% 提高到 2.54\%，AP50 从 3.84\% 提高到 8.01\%，也超过同型纯图像精修的 1.91\% 和 5.56\%。}
\newcommand{\MainResult}{在表\ref{tab:main}中，雷达补检在强模糊验证集上的 mAP50--95 为 2.54\%，比纯相机、同型纯图像精修及直接合并雷达框分别高 1.16、0.63 和 1.45 个百分点。与纯相机相比，AP50 由 3.84\% 提高至 8.01\%，AR100 由 3.17\% 提高至 11.28\%。表\ref{tab:seeds}和图\ref{fig:diagnostics}给出三个种子的结果。三次运行中，雷达补检的 mAP50--95 均高于纯相机和同型纯图像精修，显示收益方向在当前固定划分上保持一致。}
\newcommand{\DiagnosticResult}{表\ref{tab:classes}显示，相对纯相机，汽车 AP 由 3.77\% 升至 6.52\%，行人 AP 由 0.39\% 升至 1.10\%，自行车 AP 仍为零。整体改善主要来自汽车和行人。移除约一半验证帧的雷达候选后，同一补检器的 mAP50--95 从 2.54\% 降至 1.93\%，AP50 从 8.01\% 降至 5.73\%。在视觉输入保持相同的情况下，这一变化反映补检模型依赖雷达候选来恢复部分检测。}
\newcommand{\ConclusionResult}{在人工高斯强模糊下，基于雷达的独立候选经图像精修后改善了道路目标的二维检测。相对冻结的纯相机基线，三次运行的平均 mAP50--95 提高 1.16 个百分点，AP50 提高 4.17 个百分点。与同型纯图像精修相比，平均 mAP50--95 高 0.63 个百分点。部分雷达缺失时收益减小。}
